\subsection{柱投影的全层提升：边缘化一致性、pro-\texorpdfstring{$\mathrm{PerDatum}$}{PerDatum} 坐标与函子自然性}\label{subsec:op-algebra-cylinder-properdatum}
本小节不引入新的动力学假设，而是在附录 \ref{app:op-algebra} 已给出的“投影作为概念”口径下，
把长度 \(m\) 的柱投影族 \(\{P_a\}_{a\in\{0,1\}^m}\) 组织为一个全层（\(m\to\infty\)）的一致坐标系统，
并指出其对共轭态射的自然性；该结构为正文多处“柱事件/前缀/可见 \(\sigma\)-代数”的形式用法提供一个与实现无关的统一接口。

\paragraph{设置}
沿用附录开头的记号：在一个（可审计的）有限维 realization 中，取 Hilbert 空间 \(\mathcal{H}\)、幺正算子 \(U\) 与投影 \(P\)，并令
\(
P_t:=U^{-t}PU^t
\)。
对任意长度 \(m\) 的字 \(a=a_0\cdots a_{m-1}\in\{0,1\}^m\)，柱投影定义为
\[
P_a=\prod_{j=0}^{m-1}Q_j(a_j),
\qquad
Q_j(1)=P_j,\quad Q_j(0)=I-P_j
\]
（见附录 \ref{app:op-algebra} 开头）。
令 \(\tau\) 为 \(\mathcal{B}(\mathcal{H})\) 上的归一化迹 \(\tau(X):=\Tr(X)/\dim\mathcal{H}\)，并记柱概率
\[
p_m(a):=\tau(P_a)\in[0,1].
\]

\begin{proposition}[柱概率的边缘化一致性]\label{prop:op-algebra-cylinder-marginalization}
\leanverified{paper\_op\_algebra\_cylinder\_marginalization}
对任意 \(m\ge 1\) 与任意 \(a\in\{0,1\}^m\)，有精确恒等式
\[
p_m(a)=p_{m+1}(a0)+p_{m+1}(a1).
\]
因而 \(\{p_m\}_{m\ge 1}\) 形成一个边缘化一致的投影系统（Kolmogorov 一致性），
并确定一个（以柱坐标为系数的）极限概率对象。
\end{proposition}

\begin{proof}
由定义
\(
P_{a0}=P_a\,Q_m(0)
\)
与
\(
P_{a1}=P_a\,Q_m(1)
\)，
而 \(Q_m(0)+Q_m(1)=I\)，故
\(
P_{a0}+P_{a1}=P_a
\)。
两边取迹并用 \(\tau\) 的线性性即得结论。证毕。
\end{proof}

\begin{remark}[pro-\texorpdfstring{$\mathrm{PerDatum}$}{PerDatum} 的坐标化理解]\label{rem:op-algebra-properdatum-coordinate}
在定义 \ref{def:perdatum} 的语言中，
每个柱投影 \(P_a\) 给出一个周期数据对象 \((\mathcal{B}(\mathcal{H}),\tau,P_a)\)；
命题 \ref{prop:op-algebra-cylinder-marginalization} 的边缘化一致性表明：这些对象的“可观测坐标”
\(\tau(P_a)\) 在跨深度投影下严格相容，从而可被视为一类 pro-\(\mathrm{PerDatum}\) 坐标系统
（即以柱事件族作为坐标、以边缘化作为结构映射的逆系统）。
本文后续仅使用该坐标系统的可审计一致性，不依赖对 pro-范畴的额外形式化。
\end{remark}

\begin{proposition}[共轭态射下的自然性]\label{prop:op-algebra-cylinder-naturality}
设 \((\mathcal{H},U,P,\tau)\) 与 \((\mathcal{H}',U',P',\tau')\) 为两套上述数据，
且存在幺正同构 \(W:\mathcal{H}\to\mathcal{H}'\) 使
\(
WU=U'W
\)、
\(
WP=P'W
\)，
并且 \(\tau'(W X W^{-1})=\tau(X)\) 对所有 \(X\) 成立。
则对任意 \(m\) 与任意 \(a\in\{0,1\}^m\) 有
\[
W P_a W^{-1}=P'_a,
\qquad
p_m(a)=p'_m(a).
\]
因此柱坐标系统 \(\{p_m\}_m\) 对共轭实现不变。
\end{proposition}

\begin{proof}
由 \(WU=U'W\) 得 \(W P_t W^{-1}=P'_t\)（对一切 \(t\)）。
再由柱投影 \(P_a\) 对 \(\{P_t\}\) 的多项式表达式，得 \(W P_a W^{-1}=P'_a\)。
最后用 \(\tau'(W X W^{-1})=\tau(X)\) 得 \(p_m(a)=p'_m(a)\)。证毕。
\end{proof}

\begin{corollary}[折叠推前的纤维求和律（坐标层面的形式恒等式）]\label{cor:op-algebra-fold-pushforward-fiber-sum}
对任意深度 \(m\) 的折叠映射 \(\Fold_m:\Omega_m=\{0,1\}^m\to X_m\)，
若以柱坐标给定 \(\Omega_m\) 上分布 \(p_m\)，则其推前 \(\pi_m:=(\Fold_m)_*p_m\in\Delta(X_m)\) 满足
\[
\pi_m(x)=\sum_{a\in\Omega_m:\ \Fold_m(a)=x} p_m(a).
\]
该恒等式是正文中“折叠稳定化的纤维求和”口径在概率坐标上的直接版本。
\end{corollary}
